\documentclass[11pt,a4paper]{article}
\usepackage[margin=23mm]{geometry}
\usepackage{amsmath,amssymb,amsthm,hyperref}
\newtheorem{lemma}{Lemma}
\newtheorem{proposition}[lemma]{Proposition}
\newtheorem{theorem}[lemma]{Theorem}
\newcommand{\gen}[1]{\langle #1\rangle}
\title{Involution blocks and generation in the small Ree groups}
\author{Human Compression project: a partial recovery of L039}
\date{}
\begin{document}\maketitle
This note supplies complete proofs for the rank-three construction and the
odd-quotient exclusion in Leemans--Schulte--Van Maldeghem~\cite{LSV}.
A centralizer criterion also recovers their upper bound of four, using an
explicitly retained structural result. The rank-four exclusion remains
unrecovered here. We isolate a difficulty in its printed proof; neither that
difficulty nor the results below refute the paper's principal theorem.

\section{Three elementary criteria}
\begin{lemma}[Odd quotients]
If $N\trianglelefteq G$ and $G/N$ has odd order, every involution of $G$
belongs to $N$. In particular, $G$ cannot be generated by involutions if
$N\ne G$.
\end{lemma}
\begin{proof}
The image of an involution has order dividing two and the odd number
$|G/N|$, hence is the identity. Every product of involutions is therefore
in $N$.
\end{proof}

\begin{lemma}[Centralizer intersections]\label{central}
Let $G$ be a nonabelian simple group. Suppose that $C_G(a)\cap C_G(b)$
is abelian whenever $a,b$ are noncommuting involutions. Then no generating
string of involutions for $G$ has length at least five.
\end{lemma}
\begin{proof}
Suppose $(r_0,\ldots,r_{d-1})$ is such a string with $d\ge5$.
If $r_0$ commuted with $r_1$, it would commute with every generator,
contradicting the trivial center of $G$. Similarly $r_{d-2}$ and
$r_{d-1}$ do not commute. String commutation puts both of the latter
in $C_G(r_0)\cap C_G(r_1)$, an abelian group, a contradiction.
This argument needs no intersection property.
\end{proof}

\begin{lemma}[Coprime adjacent orders]\label{IP}
Let $G=\gen{a,b,c}$, where $a,b,c$ are distinct involutions,
$ac=ca$, and $|ab|=m$, $|bc|=n$ are coprime odd integers greater than
one. Suppose $G$ is larger than each of its three two-generator subgroups.
Then $(a,b,c)$ has the intersection property.
\end{lemma}
\begin{proof}
The adjacent subgroups are dihedral of orders $2m$ and $2n$. Their
intersection contains $\gen b$ and has order dividing $2\gcd(m,n)=2$,
so equals $\gen b$. The third parabolic is $\gen{a,c}\cong C_2^2$.
Its intersections with the adjacent parabolics have order dividing
$\gcd(4,2m)=\gcd(4,2n)=2$ and contain the prescribed common endpoint.
An omitted singleton generator cannot lie in its opposite parabolic,
since then that proper subgroup would contain every generator. Distinct
singleton subgroups intersect trivially. Together with the containment,
empty-set and whole-group cases, these observations give the assertion
for every pair of generator subsets.
\end{proof}

\section{A partition by fixed blocks}
Let the finite set $\Omega$ carry a Steiner $2$-design whose blocks all have the same
even size $k$. Suppose a group acts on this design and its involutions
are in bijection with the blocks, an involution $u$ corresponding to its
full fixed set $B_u$. Thus distinct involutions have at most one common
fixed point.

\begin{lemma}\label{partition}
Fix an involution $t$. The blocks of the involutions in $C_G(t)$ are
pairwise disjoint. If that centralizer has exactly $|\Omega|/k$
involutions, those blocks partition $\Omega$.
\end{lemma}
\begin{proof}
For $u\in C_G(t)\setminus\{t\}$, the set $B_t$ is invariant under $u$.
The number of its points fixed by $u$ is even, since $k$ is even and
all nonfixed orbits have size two. That number is at most one, so is
zero. Hence $B_t\cap B_u=\varnothing$.

Now take distinct $u,v\in C_G(t)\setminus\{t\}$. If $p$ belonged to
$B_u\cap B_v$, then $p\notin B_t$, so $p^t\ne p$. Commutation makes
$p^t$ another common fixed point of $u$ and $v$, contrary to the
Steiner condition. Thus all the blocks are disjoint. Under the numerical
hypothesis, their union has $(|\Omega|/k)k=|\Omega|$ points.
\end{proof}

The small Ree group $G=R(q)$, $q=3^{2e+1}\ge27$, acts on its unital,
a Steiner $2$-design with $|\Omega|=q^3+1$ and $k=q+1$. Its involution
centralizer is $\gen t\times\operatorname{PSL}_2(q)$. We retain these
classical structural facts, including the block--involution bijection.
The usual involution class calculation in $\operatorname{PSL}_2(q)$,
where $q\equiv3\pmod4$, gives
\[
 \frac{|\operatorname{PSL}_2(q)|}{q+1}=\frac{q(q-1)}2
\]
involutions. Consequently $C_G(t)$ has $1+q(q-1)$ involutions, counting
$t$ and the two lifts of each involution in the second factor. Since
\[
 (q^2-q+1)(q+1)=q^3+1,
\]
Lemma~\ref{partition} applies. In particular each point outside $B_t$
is fixed by exactly one involution different from $t$ that commutes
with $t$. The partition conclusion already appears in
Pierro~\cite[Lemma 2.2 and the following paragraph]{Pierro}; the proof
above records the direct fixed-point argument. It is not a new result
claimed by this project.

\section{The rank-three marking}
We retain the standard subgroup list for $R(q)$ used in Section 2.2
of~\cite{LSV}. Put
\[
 \alpha=q+1-\sqrt{3q},\qquad \beta=q+1+\sqrt{3q}.
\]
The two nonsplit cyclic Hall subgroups have orders $\alpha,\beta$ and
normalizers of the form $C_\alpha:C_6$ and $C_\beta:C_6$.
The other maximal subgroups are point stabilizers $U:C_{q-1}$,
involution centralizers $C_2\times\operatorname{PSL}_2(q)$, four-group
normalizers of order $6(q+1)$, and proper subfield Ree groups.
We also retain the element-order list: besides $2,3,6,9$, nonidentity
orders divide $q-1$, $(q+1)/2$, $\alpha$, or $\beta$.
These facts are summarized independently in~\cite[Section 2, Tables 1--3]{Pierro}.

\begin{proposition}\label{generation}
For every $q=3^{2e+1}\ge27$, the group $R(q)$ has a string C-group
marking $(t,u,v)$ with $|tu|=\beta$ and $|uv|\in\{3,9\}$.
\end{proposition}
\begin{proof}
Let $T=\gen x$ have order $\beta$ and let $M=N_G(T)=T:C_6$.
The involution of its complement inverts $T$. Thus its involutions
generate a dihedral subgroup $D$ of order $2\beta$. Choose involutions
$t,u\in D$ with $tu=x$.

We first verify that $M$ is the unique maximal subgroup containing
$x$. We have
\[
 \alpha\beta=q^2-q+1,
 \quad\gcd(\alpha,\beta)=\gcd(\beta,q-1)
 =\gcd(\beta,q+1)=\gcd(\beta,6)=1.
\]
For the first coprimality, a common divisor divides $2\sqrt{3q}$ and
is prime to six, so is one. For the other two, reduce $q^2-q+1$
modulo $q-1$ and $q+1$; the remainders are respectively $1$ and $3$.
Hence all the other listed maximal-subgroup types except the subfield
groups cannot contain an element of order $\beta$. A proper subfield has parameter $q_0$ with $q\ge q_0^3$.
Every element order in $R(q_0)$ is at most
$\max(9,q_0+\sqrt{3q_0}+1)$, which is smaller than $\beta$; this
also covers $q_0=3$. Finally, in any normalizer $T':C_6$ with
$|T'|=\beta$, the image of $x$ in the quotient is trivial, since its
order is prime to six. Thus $T'=\gen x=T$, and that normalizer is $M$.

The blocks $B_t$ and $B_u$ are disjoint. Otherwise $D$ fixes their
common point, whereas a point stabilizer of order $q^3(q-1)$ has
no element of order $\beta$. Choose $p\in B_u$. The partition lemma
gives a unique involution $v\ne t$ in $C_G(t)$ fixing $p$.
It also gives $v\ne u$, since $t$ and $u$ do not commute. The product
$uv$ is a nonidentity element in the normal $3$-subgroup of $G_p$:
both involutions have the same nontrivial image in its cyclic quotient
$C_{q-1}$. That $3$-subgroup has exponent nine by the retained order
list, so $|uv|=3$ or $9$.

Every involution in $M=T:C_6$ lies in the dihedral preimage of the
unique order-two subgroup of $C_6$. The product of any two such
involutions lies in $T$ and has order dividing $\beta$. Since
$|uv|\in\{3,9\}$, it follows that $v\notin M$. The unique-maximal
property of $x=tu$ now proves $\gen{t,u,v}=G$.
The adjacent orders are coprime odd numbers and $t$ commutes with
$v$. Simplicity excludes equality of $G$ with any of the dihedral
or Klein-four parabolics. Lemma~\ref{IP} proves all intersections.
\end{proof}

This recovers source Lemma 15 with its missing choice of $v$ justified.
In its last generation sentence the printed $\rho_0\notin M$ cannot
hold, since $\rho_0$ was chosen in $M$; the argument requires
$\rho_2\notin M$, which has been proved above.

\section{The upper bound and the remaining obstruction}
The classical description $\operatorname{Aut}(R(q))/R(q)\cong C_{2e+1}$
shows by the odd-quotient lemma that no proper almost simple extension
is generated by involutions. Likewise $R(3)\cong\operatorname{PSL}_2(8):C_3$
is not generated by involutions. This recovers source Lemma 2 and the
stated nonsimple exception.

For completeness about the rank bound, we retain Pierro's structural
intersection result~\cite[Lemma 3.13, noncommuting case]{Pierro}:
for noncommuting involutions $a,b$ of a simple small Ree group,
$C_G(a)\cap C_G(b)$ is elementary abelian of characteristic three,
or is a subgroup of a Klein four-group. It is therefore abelian.
Lemma~\ref{central} immediately excludes ranks at least five, recovering
source Lemma 3. The structural intersection theorem remains an external
input whose cost must be included in any subsequent measurement.

The rank-four exclusion in source Lemmas 9--14 remains a separate task.
There is a precise issue in the last paragraph of the proof of Lemma 14.
It identifies
\[
 N=N_G(D)=(E_4\times D):C_3,\qquad H=D:C_3,
\]
and says that $H$ is normal in $N$. Under the stated four-group action,
the order-three factor permutes the three nonidentity elements of $E_4$;
see also~\cite[Section 2.2.4]{Pierro}. Since $D\trianglelefteq N$, this gives
\[
 N/D\cong A_4,\qquad H/D\cong C_3.
\]
The latter is not normal in $A_4$: its four order-three subgroups are
conjugate and distinct. Normality passes to quotients, so the displayed
normality assertion cannot hold with these simultaneous identifications.
In particular its use to force the same $H$ into two overgroups needs a
different argument. This elementary obstruction does not prove that
rank-four Ree polytopes exist, and we make no such claim.

The present note proves only the stated rank-three, outer-extension,
nonsimple and rank-at-most-four conclusions, with all borrowed inputs
identified. The full principal-result contract of L039 is still pending.
There is no new conjecture answer, Lean formalization or compression factor.

\begin{thebibliography}{9}
\bibitem{LSV} D. Leemans, E. Schulte and H. Van Maldeghem,
\emph{Groups of Ree type in characteristic 3 acting on polytopes},
retained arXiv:1501.01125v1, eighteen pages, body dated 10 May 2018.
\bibitem{Pierro} E. Pierro, \emph{The M\"obius function of the small Ree groups},
Australas. J. Combin. \textbf{66}(2) (2016), 142--176,
\url{https://ajc.maths.uq.edu.au/pdf/66/ajc_v66_p142.pdf}.
\end{thebibliography}
\end{document}
